La industria de las telecomunicaciones enfrenta uno de los entornos más competitivos y saturados del mercado global. En este contexto, la retención de clientes (generalmente referida como mitigación de \textit{Churn}) ha superado la adquisición de clientes como el principal impulsor de la rentabilidad a largo plazo. 

Las infraestructuras de Machine Learning (ML) se han posicionado fuertemente como herramientas predictivas para anticipar qué usuarios tienen la mayor propensión a cancelar sus servicios. Sin embargo, la adopción operativa de estos algoritmos a menudo fracasa en la última milla por dos razones críticas: 1) la métrica de éxito de ML (usualmente Precisión o AUC) no siempre se alinea con la rentabilidad comercial real (ROI), y 2) los resultados se entregan como cajas negras probabilísticas que no brindan explicaciones accionables a los agentes humanos.

Este documento propone arquitecturar y evaluar de forma empírica una solución integral de toma de decisiones dirigida al valor de negocio (\textit{Profit-driven}). Mediante el cruce de modelos ensamblados (XGBoost), Inteligencia Artificial Explicable (XAI) basado en valores SHAP, y tableros interactivos de Inteligencia de Negocios (BI), buscamos responder a la siguiente pregunta: ¿Cómo la integración visual de explicabilidad algorítmica y matrices de coste-beneficio empodera y rentabiliza las campañas de retención en telecomunicaciones?
